% WITHDRAWN before promotion: reverse Fisher4.3(11) has a factor3 omitted by OCR.
% The claimed arbitrary-nine exclusion below is unsupported.
% Superseded by corrected B_height_two_scalar_residual_gate.tex.
\subsection{Scalar residual representations do not survive the remaining test}

\begin{lemma}[The five Hesse valuation patterns at absolute ramification two]
\label{lem:B-RD37-patterns}
\status{Candidate; complete independent audit requested.}
Let $K$ have algebraically closed residue field of characteristic three,
absolute ramification index two, and normalized valuation $v$.
Fix a primitive cube root $\zeta\in K$. For
\[
 E_t:\quad y^2+3txy+(t^3-1)y=x^3,\qquad t^3\ne1,
                                                        \tag{B-RD37.1}
\]
suppose $E_t$ is potentially supersingular. The ordered vector
$\boldsymbol a(t)=(v(t-1),v(t-\zeta),v(t-\zeta^2))$ has one
of the following five forms; the indicated permutations are allowed.
\[
\begin{array}{c|c|c}
\text{form}&\boldsymbol a(t)&\text{tame special-inertia order}\\\hline
\mathrm A&(-1,-1,-1)&4\\
\mathrm B&(0,0,0)&2\\
\mathrm C&(1,1,1)&4\\
\mathrm D&\text{a permutation of }(2,1,1)&2\\
\mathrm E&\text{a permutation of }(3,1,1)&4
\end{array}                                             \tag{B-RD37.2}
\]
The two sets of residue vectors modulo three in the last-column
classes two and four are disjoint, and remain disjoint after negating
either vector and permuting its coordinates.
\end{lemma}
\begin{proof}
The exact invariants are
\[
 c_4=9t(t^3+8),\qquad \Delta(E_t)=27(t^3-1)^3,\qquad
 j(E_t)=\frac{27t^3(t^3+8)^3}{(t^3-1)^3}.
                                                        \tag{B-RD37.3}
\]
If $v(t)=-a<0$, then $v(j)=6-3a$. Potential supersingularity
requires $v(j)>0$, so $a=1$, giving form A. If $t$ is integral
and its residue is not one, all three entries are zero, giving B.
This includes $t=0$, with the usual convention $v(0)=+\infty$.

It remains that $\overline t=1$. The three roots of unity have
pairwise difference valuation one, since
$(\zeta-1)^2=-3\zeta$. Therefore every entry is at least one,
and at most one is greater than one. Write the vector, after
permutation, as $(r,1,1)$ with $r\ge1$. Then
\[
 v(t^3-1)=r+2,
 \qquad t^3+8=(t^3-1)+9,
 \qquad v(9)=4.
\]
For $r=1$ the numerator factor also has valuation three and
$v(j)=6$. For $r=2$ its valuation is at least four and
$v(j)\ge6$, including $j=0$. For $r\ge3$ its valuation is
exactly four, and $v(j)=12-3r$; positivity permits only $r=3$.
This proves the list.

Put $S=\sum a_i(t)$. The displayed discriminant has valuation
$\delta=6+3S$, respectively $-3,6,15,18,21$ in the five forms.
Proposition~\ref{prop:B-RD36} gives the last column; negative raw
discriminant valuations are permitted since the same proposition
applies to any Weierstrass equation. Modulo three the order-two
vectors are $(0,0,0)$ and permutations of $(2,1,1)$. The
order-four vectors are $(2,2,2)$, $(1,1,1)$ and permutations of
$(0,1,1)$. Negating the second order-two form gives permutations
of $(1,2,2)$, which occur in neither order-four list. The asserted
disjointness follows, also for arbitrary coordinate permutations.
\end{proof}

\begin{proposition}[B-RD37: scalar residual representations are excluded]
\label{prop:B-RD37}
\status{Candidate; complete independent audit requested.}
Let $K$ be as in the preceding lemma. Let $E,E'/K$ be potentially
supersingular elliptic curves with $E[9]\simeq E'[9]$.
If their common residual representation on three-torsion is scalar,
then their finite special-inertia characters agree on all inertia
elements of order prime to three.

Consequently, in the remaining level-$27$ comparison of
Theorem~\ref{thm:B-RD34}, the common residual image is non-scalar.
After a simultaneous quadratic twist it is either the nontrivial
unipotent group $C_3$ or the irreducible tame cyclic group $C_4$.
This excludes all scalar residual representations from the remaining
joint-inertia orders $18$ and $36$.
\end{proposition}
\begin{proof}
First suppose the three-torsion is rational. The full-level moduli
description of $Y(3)$ and its elliptic family, in
\citet[Section~2, p.~4; Lemma~4.2(i), p.~8]{Fisher2015Nine},
allows us to write $E$ as (B-RD37.1), with parameter $t_E\in K$.
An arbitrary isomorphism of nine-torsion has a Weil-pairing multiplier
in $(\mathbf Z/9)^\times$. Every such multiplier is plus or minus
a square. Composing with a scalar automorphism of the target torsion
therefore makes the isomorphism either symplectic or antisymplectic.
These are exactly the direct and reverse nine-congruences in
\citet[Introduction, p.~1; Section~1, first paragraph, p.~2]
{Fisher2015Nine}.

The precise original equations of
\citet[Theorem~4.3, equations~(10)--(11), pp.~9--10]
{Fisher2015Nine}, with its parameter $\Delta=1$, give a parameter
$t\in K$ representing $E'$ by the same family (B-RD37.1).
For every $\theta\in\{1,\zeta,\zeta^2\}$ they impose respectively
\[
 \frac{t-\theta}{t_E-\theta}\in(K^\times)^3
       \quad\text{in the direct case},
 \qquad
 (t-\theta^{-1})(t_E-\theta)\in(K^\times)^3
       \quad\text{in the reverse case}.
                                                        \tag{B-RD37.4}
\]
These are necessary equations for the actual nine-congruence, not
an inference from equality of the mod-three representations.
In particular the two valuation vectors agree modulo three in the
direct case. In the reverse case they are negatives modulo three,
with the permutation $\theta\mapsto\theta^{-1}$.
Lemma~\ref{lem:B-RD37-patterns} excludes a pair in different tame
columns in either case. Their tame quotient orders are equal, so
Lemma~\ref{lem:B-RD32-norm} identifies their prime-to-three trace
functions on the original inertia group, not just their abstract
character tables.

For a scalar residual representation, write it as
$\varepsilon\operatorname{id}$ with
$\varepsilon:G_K\to\mathbf F_3^\times=\{\pm1\}$.
Use the same quadratic character to twist both elliptic curves.
Their three-torsion becomes rational and the nine-torsion
isomorphism persists. Potential supersingularity also persists.
Over a common good-reduction field the new descent action is the
old action multiplied by the scalar automorphism
$[\varepsilon(\sigma)]$. Hence both special-inertia characters
are multiplied by the same function $\varepsilon$; equality after
twisting is equivalent to equality before twisting. This proves
the first assertion.

For completeness, after the completed maximal unramified extension
the entire Galois group is inertia. The common residual image lies
in $\operatorname{SL}_2(\mathbf F_3)$ because $\mu_3\subset K$.
Its wild subgroup is a normal three-subgroup and its quotient is
cyclic of order prime to three. If the wild subgroup is nontrivial,
it has order three and its normalizer in
$\operatorname{SL}_2(\mathbf F_3)$ is $C_6$; the image is $C_3$
or $C_6$. A simultaneous quadratic twist removes the central
minus sign in the latter case. If the wild subgroup is trivial,
the image is cyclic of order $1,2$ or $4$. The first two are
scalar and have just been excluded. An order-four generator has
irreducible characteristic polynomial $X^2+1$ over $\mathbf F_3$.
These elementary subgroup assertions are also exhaustively checked
by the retained finite certificate. Thus precisely the two stated
residual forms remain as possible tests; their existence as
opposite-character level-$27$ pairs is not asserted.
\end{proof}

\begin{remark}[Compatibility with the known level-nine ambiguity]
\label{rem:B-RD37-known-example}
The pair in Proposition~\ref{prop:B-RD29-p3-level-two} consists of
a potentially supersingular elliptic group and a Tate group with
good rank-one graded pieces. Moreover its common three-torsion is
non-scalar. Indeed its Kummer extension parameter is
$q=2^5 3^6$, whose cube class is nontrivial even after completed
unramified base change. If $2$ were a cube, so would be $-2$.
A putative unit $x$ with $x^3=-2$ has residue one; writing
$x=1+w$ gives $v(w)\ge1$ and
\[
 v\bigl((1+w)^3-1\bigr)
   =v(3w+3w^2+w^3)\ge3,
\]
whereas $v(-2-1)=v(-3)=2$. This is impossible.
Thus the example has nontrivial unipotent residual image $C_3$,
and supplies no exception to Proposition~\ref{prop:B-RD37}.
\end{remark}
